\documentclass[11pt,a4paper]{article}
\usepackage{fontspec,amsmath,amssymb,amsthm,microtype}
\setmainfont{lmroman10-regular.otf}[BoldFont=lmroman10-bold.otf,ItalicFont=lmroman10-italic.otf,BoldItalicFont=lmroman10-bolditalic.otf]
\usepackage{xcolor}
\usepackage[most]{tcolorbox}
\usepackage{mathrsfs}
\providecommand{\lt}{<}
\providecommand{\gt}{>}
\setmonofont{DejaVuSansMono.ttf}[BoldFont=DejaVuSansMono-Bold.ttf,ItalicFont=DejaVuSansMono-Oblique.ttf,BoldItalicFont=DejaVuSansMono-BoldOblique.ttf,Scale=0.86]
\newfontfamily\leanfallback{STIXTwoMath-Regular.otf}[Scale=0.86]
\definecolor{leanrule}{gray}{0.62}
\definecolor{leanback}{gray}{0.962}
\definecolor{leantitle}{gray}{0.22}
\definecolor{leanmuted}{gray}{0.38}
\newtcolorbox{leanbox}[1]{enhanced,breakable,colback=leanback,colframe=leanrule,
  boxrule=0.4pt,arc=0.6mm,left=2.4mm,right=2.4mm,top=1.8mm,bottom=1.8mm,
  before skip=0.8em,after skip=1em,fonttitle=\small\bfseries,coltitle=leantitle,
  colbacktitle=leanback,titlerule=0.3pt,title={#1}}
\newlength{\leancharwidth}
\newenvironment{leancode}{\par\smallskip\begingroup\ttfamily\footnotesize\raggedright
  \settowidth{\leancharwidth}{M}\setlength{\parskip}{0pt}\setlength{\parindent}{0pt}
  \frenchspacing}{\par\endgroup\smallskip}
\newcommand{\leanline}[2]{\hangindent=\dimexpr#1\leancharwidth+4\leancharwidth\relax
  \hangafter=1\noindent\hspace*{#1\leancharwidth}#2\par}
\newcommand{\leansource}[1]{\par{\footnotesize\color{leanmuted}#1}\par}
\newcommand{\leanlabel}[1]{\par\smallskip\noindent{\small\bfseries #1}\par}
\usepackage[margin=25mm]{geometry}
\usepackage[unicode,hidelinks]{hyperref}
\hypersetup{pdfauthor={LangevinLean Project}}
\setlength{\emergencystretch}{3em}
\setlength{\parskip}{0.4em}
\setlength{\parindent}{0pt}
\allowdisplaybreaks[2]
\raggedbottom
\title{\texorpdfstring{Transformed contraction for a strongly convex smooth potential}{Transformed contraction for a strongly convex smooth potential}}
\author{LangevinLean Project}
\date{}
\begin{document}\maketitle
Draft mathematical exposition; these notes may contain mistakes.


A sector inequality for the gradient difference gives contraction without choosing a time-dependent averaged Hessian. Assume the potential is differentiable and strongly convex and has Lipschitz gradient. Its two-point strong-smooth inequality yields transformed dissipativity directly for the difference equation after common forcing cancels.



Integration gives the pathwise position estimate with factor \(\sqrt2e^{-mt/\gamma}\) in the stated friction regime. Construction of stochastic integral solutions and the infimum over initial probability couplings are used separately to obtain a Wasserstein estimate. No continuous Hessian is required for this argument.



Earlier results used here: \cite[Section~3.10]{LangevinLean2026ULASC}.

\section{\texorpdfstring{The energy derivative and the gradient sector inequality}{The energy derivative and the gradient sector inequality}}


For velocity difference \(v\), position difference \(x\), and force difference \(g\), define \(D=2\langle v+\gamma x,-g\rangle+2\langle-v,\gamma v+g\rangle\). The strong-smooth interpolation inequality is \(mM\|x\|^2+\|g\|^2\le(m+M)\langle g,x\rangle\). Subtract the linear part \(mx\), expand its squared norm, and rearrange to obtain \(\|g-mx\|^2\le(M-m)\langle x,g-mx\rangle\). The proof performs this expansion with the Hilbert norm-square and inner-product identities, then uses elementary real inequalities. This sector estimate contains the joint lower–upper curvature information needed for contraction.



Write the displacement as \(q=x-y\) and the force difference as \(d=\nabla f(x)-\nabla f(y)\). Strong convexity and smoothness place \((q,d)\) in a sector narrower than a bare Lipschitz ball. This relation implies dissipation of the transformed energy.

\begin{leanbox}{Lean: \texttt{forceEnergyDerivative}, \texttt{interpolation\_\allowbreak{}sector}}
\leansource{\texttt{forceEnergyDerivative} \quad KineticLangevin/\allowbreak{}PotentialContraction.lean, lines 23–24}
\begin{leancode}
\leanline{0}{def\ forceEnergyDerivative\ (γ\ :\ ℝ)\ (v\ x\ force\ :\ E)\ :\ ℝ\ :=}
\leanline{2}{2\ *\ {\leanfallback ⟪}v\ +\ γ\ •\ x,\ -force{\leanfallback ⟫}\ +\ 2\ *\ {\leanfallback ⟪}-v,\ γ\ •\ v\ +\ force{\leanfallback ⟫}}
\end{leancode}
\leansource{\texttt{interpolation\_\allowbreak{}sector} \quad KineticLangevin/\allowbreak{}PotentialContraction.lean, lines 28–34}
\begin{leancode}
\leanline{0}{theorem\ interpolation\_sector\ (m\ M\ :\ ℝ)\ (x\ force\ :\ E)}
\leanline{4}{(hi\ :\ m\ *\ M\ *\ ‖x‖\ \textasciicircum{}\ 2\ +\ ‖force‖\ \textasciicircum{}\ 2\ ≤\ (m\ +\ M)\ *\ {\leanfallback ⟪}force,\ x{\leanfallback ⟫})\ :}
\leanline{4}{‖force\ -\ m\ •\ x‖\ \textasciicircum{}\ 2\ ≤\ (M\ -\ m)\ *\ {\leanfallback ⟪}x,\ force\ -\ m\ •\ x{\leanfallback ⟫}}
\end{leancode}
\leanlabel{How the Lean proof follows the argument}
\texttt{force\allowbreak{}Energy\allowbreak{}Derivative} names the mathematical object defined above. Its defining expression appears in the adjacent statement; subsequent identities may unfold this definition directly.




\texttt{interpolation\_\allowbreak{}sector} gives the formal counterpart of this step. Its remaining equalities are discharged by the simplifications displayed in the proof.
\end{leanbox}

\section{\texorpdfstring{Complete a square to prove dissipativity}{Complete a square to prove dissipativity}}


Put \(k=g-mx\). Expanding the transformed energy yields

\[\gamma D+2m\mathcal E
=2\bigl[-(\gamma^2-2m)\|v\|^2-2\gamma\langle v,k\rangle-\gamma^2\langle x,k\rangle\bigr].\]

When \(M=m\), the sector inequality forces \(k=0\); friction gives \(\gamma^2\ge2m\), proving nonpositivity. When \(M>m\), combine the sector estimate with \(\|\gamma k+(M-m)v\|^2\ge0\) and \((M-m)(\gamma^2-m-M)\|v\|^2\ge0\). These imply \(\gamma D+2m\mathcal E\le0\). The argument handles the equal-curvature case separately, expands the square in the strict case, and combines the scalar inequalities. Positive \(\gamma\) then permits division, giving \(D\le-2m\mathcal E/\gamma\).



The potentially positive mixed term is paired with the negative position term and rewritten as a square plus a curvature contribution. The sector inequality controls the square. Using this relation retains the strong-convexity coefficient in the differential inequality.

\begin{leanbox}{Lean: \texttt{force\_\allowbreak{}dissipativity}}
\leansource{\texttt{force\_\allowbreak{}dissipativity} \quad KineticLangevin/\allowbreak{}PotentialContraction.lean, lines 38–83}
\begin{leancode}
\leanline{0}{theorem\ force\_dissipativity\ (γ\ m\ M\ :\ ℝ)\ (hγ\ :\ 0\ <\ γ)\ (hmM\ :\ m\ ≤\ M)}
\leanline{4}{(hfriction\ :\ m\ +\ M\ ≤\ γ\ \textasciicircum{}\ 2)\ (v\ x\ force\ :\ E)}
\leanline{4}{(hsector\ :\ ‖force\ -\ m\ •\ x‖\ \textasciicircum{}\ 2\ ≤\ (M\ -\ m)\ *\ {\leanfallback ⟪}x,\ force\ -\ m\ •\ x{\leanfallback ⟫})\ :}
\leanline{4}{forceEnergyDerivative\ γ\ v\ x\ force\ ≤}
\leanline{6}{(-2\ *\ m\ /\ γ)\ *\ transformEnergy\ γ\ 0\ (v,\ x)}
\end{leancode}
\leanlabel{How the Lean proof follows the argument}
\texttt{force\_\allowbreak{}dissipativity} gives the formal counterpart of this step. The main intermediate facts are \texttt{k}, \texttt{hk}, \texttt{hid}, \texttt{hb}, \texttt{hdiv} in that order. Named proof applications include \texttt{norm\_\allowbreak{}eq\_\allowbreak{}zero.\allowbreak{}mp}, \texttt{mul\_\allowbreak{}nonneg}; the adjacent proof shows their hypotheses and order of use.
\end{leanbox}

\section{\texorpdfstring{The sector condition from the potential}{The sector condition from the potential}}


For \(g=\nabla f(x)-\nabla f(y)\) and displacement \(x-y\), the strong-smooth interpolation inequality follows from differentiability, strong convexity and gradient Lipschitzness. the strong-smooth force dissipativity inequality applies that theorem, converts it to the sector condition, and applies the completed-square estimate to that sector inequality. Consequently no Hessian, sector bound, or contraction conclusion is assumed in this potential theorem. The explicit hypotheses retain \(m\le M\), \(\gamma>0\), and \(\gamma^2\ge m+M\).



One may apply the strong-smooth interpolation inequality directly to \(f\). No continuous Hessian is required: differentiability, strong convexity, and a Lipschitz gradient already imply the needed two-point estimate. Thus this proof applies to the regularity assumptions of the first-order sampling theorem.

\begin{leanbox}{Lean: \texttt{potential\_\allowbreak{}force\_\allowbreak{}dissipativity}}
\leansource{\texttt{potential\_\allowbreak{}force\_\allowbreak{}dissipativity} \quad KineticLangevin/\allowbreak{}PotentialContraction.lean, lines 87–96}
\begin{leancode}
\leanline{0}{theorem\ potential\_force\_dissipativity\ [CompleteSpace\ E]}
\leanline{4}{(f\ :\ E\ →\ ℝ)\ (m\ γ\ :\ ℝ)\ (M\ :\ ℝ≥0)}
\leanline{4}{(hdiff\ :\ Differentiable\ ℝ\ f)\ (hsc\ :\ StrongConvexOn\ univ\ m\ f)}
\leanline{4}{(hLip\ :\ LipschitzWith\ M\ (gradient\ f))\ (hmM\ :\ m\ ≤\ (M\ :\ ℝ))}
\leanline{4}{(hγ\ :\ 0\ <\ γ)\ (hfriction\ :\ m\ +\ (M\ :\ ℝ)\ ≤\ γ\ \textasciicircum{}\ 2)\ (v\ x\ y\ :\ E)\ :}
\leanline{4}{forceEnergyDerivative\ γ\ v\ (x\ -\ y)\ (gradient\ f\ x\ -\ gradient\ f\ y)\ ≤}
\leanline{6}{(-2\ *\ m\ /\ γ)\ *\ transformEnergy\ γ\ 0\ (v,\ x\ -\ y)}
\end{leancode}
\leanlabel{How the Lean proof follows the argument}
\texttt{potential\_\allowbreak{}force\_\allowbreak{}dissipativity} gives the formal counterpart of this step. Named proof applications include \texttt{force\_\allowbreak{}dissipativity}, \texttt{interpolation\_\allowbreak{}sector}; the adjacent proof shows their hypotheses and order of use.
\end{leanbox}


Earlier results used here: \cite[Section~3.10.2]{LangevinLean2026ULASC}.

\section{\texorpdfstring{Differentiate the noise-cancelled path energy}{Differentiate the noise-cancelled path energy}}


Assume \(v'=-\gamma v-g\) and \(x'=v\). Then \((v+\gamma x)'=-g\), while \((-v)'=\gamma v+g\). The argument obtains these two derivative formulas by addition, negation and constant scalar multiplication, normalizes them with vector-space algebra, and applies the squared-norm derivative rule to each coordinate. Their sum is precisely the energy derivative. The theorem assumes differentiability of the difference paths; it does not incorrectly differentiate an individual Brownian velocity path.



For two paths driven by identical forcing, subtract their integral equations first. Their position difference has derivative equal to their velocity difference, and their velocity difference has derivative equal to the negative friction and force difference. The energy chain rule is then an ordinary deterministic calculation.

\begin{leanbox}{Lean: \texttt{hasDerivAt\_\allowbreak{}forceEnergy}}
\leansource{\texttt{hasDerivAt\_\allowbreak{}forceEnergy} \quad KineticLangevin/\allowbreak{}PotentialContraction.lean, lines 100–114}
\begin{leancode}
\leanline{0}{theorem\ hasDerivAt\_forceEnergy\ (γ\ :\ ℝ)\ (v\ x\ force\ :\ ℝ\ →\ E)\ (t\ :\ ℝ)}
\leanline{4}{(hv\ :\ HasDerivAt\ v\ (-γ\ •\ v\ t\ -\ force\ t)\ t)}
\leanline{4}{(hx\ :\ HasDerivAt\ x\ (v\ t)\ t)\ :}
\leanline{4}{HasDerivAt\ (fun\ r\ ↦\ transformEnergy\ γ\ 0\ (v\ r,\ x\ r))}
\leanline{6}{(forceEnergyDerivative\ γ\ (v\ t)\ (x\ t)\ (force\ t))\ t}
\end{leancode}
\leanlabel{How the Lean proof follows the argument}
\texttt{has\allowbreak{}Deriv\allowbreak{}At\_\allowbreak{}force\allowbreak{}Energy} gives the formal counterpart of this step. The main intermediate facts are \texttt{hy}, \texttt{hz}, \texttt{hyd}, \texttt{hzd} in that order. Its remaining equalities are discharged by the simplifications displayed in the proof.
\end{leanbox}

\section{\texorpdfstring{Integrate energy decay and take the position square root}{Integrate energy decay and take the position square root}}


For two continuous position paths \(x,y\), set \(F(t)=\mathcal E(v_t,x_t-y_t)\). The derivative theorem and potential dissipativity give \(F'\le-2mF/\gamma\). Closed-interval continuity and the derived derivative formulas allow the same Gronwall theorem as in the spectral proof, giving \(F(t)\le F(0)e^{-2mt/\gamma}\). If \(v_0=0\), then \(F(0)=\gamma^2\|x_0-y_0\|^2\). Combine with the position projection bound, cancel the positive \(\gamma^2\), and compare nonnegative squares. The exponential square identity \((e^{-mt/\gamma})^2=e^{-2mt/\gamma}\) yields \(\|x_t-y_t\|\le\sqrt2e^{-mt/\gamma}\|x_0-y_0\|\). This is a proved nonlinear pathwise estimate with the paper's correct prefactor.



After the differential inequality is established, integration gives exponential decay of the squared energy. Taking a square root halves its exponent. The equal-initial-velocity identity and the position projection then give the visible \(\sqrt2\) coefficient in the simple friction regime.

\begin{leanbox}{Lean: \texttt{potential\_\allowbreak{}synchronous\_\allowbreak{}energy\_\allowbreak{}decay}, \texttt{potential\_\allowbreak{}synchronous\_\allowbreak{}position\_\allowbreak{}decay}}
\leansource{\texttt{potential\_\allowbreak{}synchronous\_\allowbreak{}energy\_\allowbreak{}decay} \quad KineticLangevin/\allowbreak{}PotentialContraction.lean, lines 119–153}
\begin{leancode}
\leanline{0}{theorem\ potential\_synchronous\_energy\_decay\ [CompleteSpace\ E]}
\leanline{4}{(f\ :\ E\ →\ ℝ)\ (m\ γ\ :\ ℝ)\ (M\ :\ ℝ≥0)}
\leanline{4}{(hdiff\ :\ Differentiable\ ℝ\ f)\ (hsc\ :\ StrongConvexOn\ univ\ m\ f)}
\leanline{4}{(hLip\ :\ LipschitzWith\ M\ (gradient\ f))\ (hmM\ :\ m\ ≤\ (M\ :\ ℝ))}
\leanline{4}{(hγ\ :\ 0\ <\ γ)\ (hfriction\ :\ m\ +\ (M\ :\ ℝ)\ ≤\ γ\ \textasciicircum{}\ 2)}
\leanline{4}{(T\ :\ ℝ)\ (v\ x\ y\ :\ ℝ\ →\ E)}
\leanline{4}{(hvcont\ :\ ContinuousOn\ v\ (Icc\ 0\ T))}
\leanline{4}{(hxcont\ :\ ContinuousOn\ x\ (Icc\ 0\ T))\ (hycont\ :\ ContinuousOn\ y\ (Icc\ 0\ T))}
\leanline{4}{(hv\ :\ ∀\ t\ ∈\ Ico\ 0\ T,\ HasDerivAt\ v}
\leanline{6}{(-γ\ •\ v\ t\ -\ (gradient\ f\ (x\ t)\ -\ gradient\ f\ (y\ t)))\ t)}
\leanline{4}{(hx\ :\ ∀\ t\ ∈\ Ico\ 0\ T,\ HasDerivAt\ (fun\ r\ ↦\ x\ r\ -\ y\ r)\ (v\ t)\ t)}
\leanline{4}{(t\ :\ ℝ)\ (ht\ :\ t\ ∈\ Icc\ 0\ T)\ :}
\leanline{4}{transformEnergy\ γ\ 0\ (v\ t,\ x\ t\ -\ y\ t)\ ≤}
\leanline{6}{transformEnergy\ γ\ 0\ (v\ 0,\ x\ 0\ -\ y\ 0)\ *\ Real.exp\ ((-2\ *\ m\ /\ γ)\ *\ t)}
\end{leancode}
\leansource{\texttt{potential\_\allowbreak{}synchronous\_\allowbreak{}position\_\allowbreak{}decay} \quad KineticLangevin/\allowbreak{}PotentialContraction.lean, lines 158–186}
\begin{leancode}
\leanline{0}{theorem\ potential\_synchronous\_position\_decay\ [CompleteSpace\ E]}
\leanline{4}{(f\ :\ E\ →\ ℝ)\ (m\ γ\ :\ ℝ)\ (M\ :\ ℝ≥0)}
\leanline{4}{(hdiff\ :\ Differentiable\ ℝ\ f)\ (hsc\ :\ StrongConvexOn\ univ\ m\ f)}
\leanline{4}{(hLip\ :\ LipschitzWith\ M\ (gradient\ f))\ (hmM\ :\ m\ ≤\ (M\ :\ ℝ))}
\leanline{4}{(hγ\ :\ 0\ <\ γ)\ (hfriction\ :\ m\ +\ (M\ :\ ℝ)\ ≤\ γ\ \textasciicircum{}\ 2)}
\leanline{4}{(T\ :\ ℝ)\ (v\ x\ y\ :\ ℝ\ →\ E)\ (hv0\ :\ v\ 0\ =\ 0)}
\leanline{4}{(hvcont\ :\ ContinuousOn\ v\ (Icc\ 0\ T))}
\leanline{4}{(hxcont\ :\ ContinuousOn\ x\ (Icc\ 0\ T))\ (hycont\ :\ ContinuousOn\ y\ (Icc\ 0\ T))}
\leanline{4}{(hv\ :\ ∀\ t\ ∈\ Ico\ 0\ T,\ HasDerivAt\ v}
\leanline{6}{(-γ\ •\ v\ t\ -\ (gradient\ f\ (x\ t)\ -\ gradient\ f\ (y\ t)))\ t)}
\leanline{4}{(hx\ :\ ∀\ t\ ∈\ Ico\ 0\ T,\ HasDerivAt\ (fun\ r\ ↦\ x\ r\ -\ y\ r)\ (v\ t)\ t)}
\leanline{4}{(t\ :\ ℝ)\ (ht\ :\ t\ ∈\ Icc\ 0\ T)\ :}
\leanline{4}{‖x\ t\ -\ y\ t‖\ ≤\ Real.sqrt\ 2\ *\ Real.exp\ ((-m\ /\ γ)\ *\ t)\ *\ ‖x\ 0\ -\ y\ 0‖}
\end{leancode}
\leanlabel{How the Lean proof follows the argument}
\texttt{potential\_\allowbreak{}synchronous\_\allowbreak{}energy\_\allowbreak{}decay} gives the formal counterpart of this step. The main intermediate facts are \texttt{F}, \texttt{hcont}, \texttt{hder}, \texttt{hbound}, \texttt{hgr} in that order. Named proof applications include \texttt{has\allowbreak{}Deriv\allowbreak{}At\_\allowbreak{}force\allowbreak{}Energy}, \texttt{potential\_\allowbreak{}force\_\allowbreak{}dissipativity}, \texttt{hgr}; the adjacent proof shows their hypotheses and order of use.




\texttt{potential\_\allowbreak{}synchronous\_\allowbreak{}position\_\allowbreak{}decay} gives the formal counterpart of this step. The main intermediate facts are \texttt{he}, \texttt{hp}, \texttt{hs}, \texttt{hexp} in that order. Named proof applications include \texttt{hs}; the adjacent proof shows their hypotheses and order of use.
\end{leanbox}

\section{\texorpdfstring{Use right derivatives at the initial time}{Use right derivatives at the initial time}}


A finite-horizon integral solution is constrained only for nonnegative time. Its extension to negative time need not have the same derivative. The correct initial derivative is therefore the derivative within the right half-line. The product and squared-norm derivative rules remain valid in this form, so differentiating the transformed energy gives the same force-energy expression. The argument applies those within-derivative rules to both transformed coordinates, cancels the opposing friction terms, and transfers the resulting derivative to the energy expression. The right-slope form of Gronwall requires exactly these hypotheses, and gives the same exponential energy decay throughout the closed interval. This is the version used by the integral-flow theorem.



At time zero all derivatives are interpreted from within the nonnegative time interval. This is sufficient for the integrating-factor argument on a closed finite horizon. It avoids extending either trajectory or its driving noise to negative times.

\begin{leanbox}{Lean: \texttt{hasDerivWithinAt\_\allowbreak{}forceEnergy}, \texttt{potential\_\allowbreak{}synchronous\_\allowbreak{}energy\_\allowbreak{}decay\_\allowbreak{}right}}
\leansource{\texttt{hasDerivWithinAt\_\allowbreak{}forceEnergy} \quad KineticLangevin/\allowbreak{}PotentialContraction.lean, lines 190–204}
\begin{leancode}
\leanline{0}{theorem\ hasDerivWithinAt\_forceEnergy\ (γ\ :\ ℝ)\ (v\ x\ force\ :\ ℝ\ →\ E)\ (t\ :\ ℝ)}
\leanline{4}{(hv\ :\ HasDerivWithinAt\ v\ (-γ\ •\ v\ t\ -\ force\ t)\ (Ici\ t)\ t)}
\leanline{4}{(hx\ :\ HasDerivWithinAt\ x\ (v\ t)\ (Ici\ t)\ t)\ :}
\leanline{4}{HasDerivWithinAt\ (fun\ r\ ↦\ transformEnergy\ γ\ 0\ (v\ r,\ x\ r))}
\leanline{6}{(forceEnergyDerivative\ γ\ (v\ t)\ (x\ t)\ (force\ t))\ (Ici\ t)\ t}
\end{leancode}
\leansource{\texttt{potential\_\allowbreak{}synchronous\_\allowbreak{}energy\_\allowbreak{}decay\_\allowbreak{}right} \quad KineticLangevin/\allowbreak{}PotentialContraction.lean, lines 209–243}
\begin{leancode}
\leanline{0}{theorem\ potential\_synchronous\_energy\_decay\_right\ [CompleteSpace\ E]}
\leanline{4}{(f\ :\ E\ →\ ℝ)\ (m\ γ\ :\ ℝ)\ (M\ :\ ℝ≥0)}
\leanline{4}{(hdiff\ :\ Differentiable\ ℝ\ f)\ (hsc\ :\ StrongConvexOn\ univ\ m\ f)}
\leanline{4}{(hLip\ :\ LipschitzWith\ M\ (gradient\ f))\ (hmM\ :\ m\ ≤\ (M\ :\ ℝ))}
\leanline{4}{(hγ\ :\ 0\ <\ γ)\ (hfriction\ :\ m\ +\ (M\ :\ ℝ)\ ≤\ γ\ \textasciicircum{}\ 2)}
\leanline{4}{(T\ :\ ℝ)\ (v\ x\ y\ :\ ℝ\ →\ E)}
\leanline{4}{(hvcont\ :\ ContinuousOn\ v\ (Icc\ 0\ T))}
\leanline{4}{(hxcont\ :\ ContinuousOn\ x\ (Icc\ 0\ T))\ (hycont\ :\ ContinuousOn\ y\ (Icc\ 0\ T))}
\leanline{4}{(hv\ :\ ∀\ t\ ∈\ Ico\ 0\ T,\ HasDerivWithinAt\ v}
\leanline{6}{(-γ\ •\ v\ t\ -\ (gradient\ f\ (x\ t)\ -\ gradient\ f\ (y\ t)))\ (Ici\ t)\ t)}
\leanline{4}{(hx\ :\ ∀\ t\ ∈\ Ico\ 0\ T,\ HasDerivWithinAt\ (fun\ r\ ↦\ x\ r\ -\ y\ r)\ (v\ t)\ (Ici\ t)\ t)}
\leanline{4}{(t\ :\ ℝ)\ (ht\ :\ t\ ∈\ Icc\ 0\ T)\ :}
\leanline{4}{transformEnergy\ γ\ 0\ (v\ t,\ x\ t\ -\ y\ t)\ ≤}
\leanline{6}{transformEnergy\ γ\ 0\ (v\ 0,\ x\ 0\ -\ y\ 0)\ *\ Real.exp\ ((-2\ *\ m\ /\ γ)\ *\ t)}
\end{leancode}
\leanlabel{How the Lean proof follows the argument}
\texttt{has\allowbreak{}Deriv\allowbreak{}Within\allowbreak{}At\_\allowbreak{}force\allowbreak{}Energy} gives the formal counterpart of this step. The main intermediate facts are \texttt{hy}, \texttt{hz}, \texttt{hyd}, \texttt{hzd} in that order. Its remaining equalities are discharged by the simplifications displayed in the proof.




\texttt{potential\_\allowbreak{}synchronous\_\allowbreak{}energy\_\allowbreak{}decay\_\allowbreak{}right} gives the formal counterpart of this step. The main intermediate facts are \texttt{F}, \texttt{hcont}, \texttt{hder}, \texttt{hbound}, \texttt{hgr} in that order. Named proof applications include \texttt{has\allowbreak{}Deriv\allowbreak{}Within\allowbreak{}At\_\allowbreak{}force\allowbreak{}Energy}, \texttt{potential\_\allowbreak{}force\_\allowbreak{}dissipativity}, \texttt{hgr}; the adjacent proof shows their hypotheses and order of use.
\end{leanbox}

\begin{thebibliography}{ULA-LSI}
\bibitem[ULA-SC]{LangevinLean2026ULASC}
LangevinLean Project. \emph{Unadjusted Langevin under Strong Convexity}. LL-01, 2026. Draft report, content revision d34ead06987c. \url{https://odakyildiz.github.io/LangevinLean/docs/pdf/UnadjustedLangevin/FullProof.pdf}.
\end{thebibliography}

\end{document}
